\section{Finite verification}\label{sec:finite-verification}

The certificates in Section~\ref{sec:polynomial-certificates} reduce the two
product inequalities to signs of rational coefficient arrays.  We now specify
a finite subdivision of $\mathcal S\times\mathcal S$ on which those signs can
be checked.  The subdivision and all coefficient operations are exact.  The
complete verifier is printed in Appendix~\ref{app:verifiers}.

\subsection{Subdivision and coverage}

A cell is specified by two affinely independent lists of four rational vertices
$a=(a_0,a_1,a_2,a_3)$ and $b=(b_0,b_1,b_2,b_3)$ in $\mathcal S$; its domain is
$\operatorname{conv}(a)\times\operatorname{conv}(b)$.  Initially both lists
are $(S_0,S_1,S_2,S_3)$.  On a cell we substitute
\[
 p=\sum_{i=0}^3X_i a_i,\qquad q=\sum_{j=0}^3Y_j b_j,
 \qquad X,Y\in\Deltafour,
\]
into the arrays of Section~\ref{sec:polynomial-certificates}.

\begin{lemma}[Midpoint subdivision]\label{lem:subdivision}
Let $\varepsilon_0,\ldots,\varepsilon_3$ be the unit vectors in the
barycentric coordinate space.  Replace vertex $a_i$ by $(a_i+a_j)/2$, where
$i\ne j$, leaving the other vertices unchanged.  For an ordinary homogeneous coefficient indexed by
$I\in\mathcal I_r$, its contributions to the new coefficients are indexed by
\begin{equation}\label{eq:midpoint-substitution}
 I+t\varepsilon_i-t\varepsilon_j,\qquad 0\le t\le I_j,
 \qquad\text{with multiplier}\quad
 2^{-I_i-t}\binom{I_j}{t}.
\end{equation}
The same formula holds in the $Y$ index when a vertex of $b$ is replaced.
The two cells obtained by replacing vertex $i$, respectively vertex $j$,
by their midpoint cover the original cell.
\end{lemma}

\begin{proof}
Write $X'$ for the new barycentric coordinates.  The old coordinates satisfy
$X_i=X'_i/2$, $X_j=X'_j+X'_i/2$, with the others unchanged.  Expanding
$X_i^{I_i}X_j^{I_j}$ gives \eqref{eq:midpoint-substitution}.  The first child
contains precisely the barycentric points with $X_i\le X_j$: one can take
$X'_i=2X_i$ and $X'_j=X_j-X_i$.  The second child contains the points with
$X_j\le X_i$.  These two inequalities exhaust the original simplex.
Each replacement preserves affine independence, since its barycentric
change of coordinates has nonzero determinant.
\end{proof}

Here is the entire decision rule for one row of \eqref{eq:rational-rows}.
Start at depth zero.  At each cell:
\begin{enumerate}
\item Try the Type T coefficient test of Lemma~\ref{lem:type-t}.  Accept the
cell if every coefficient is nonnegative.
\item If neither vertex list contains $\unif$, try the Type P test of
Lemma~\ref{lem:type-p}, with $M=2^{32}$.  If the correction
\eqref{eq:remainder-correction} is undefined, report failure immediately.
Otherwise accept the cell if its coefficient test succeeds.
\item If neither test has accepted the cell and its depth is 24, report
failure.  Otherwise select an edge as follows.  For $y\in\{0,1\}$, put
$c=a$ when $y=0$ and $c=b$ when $y=1$.  Among all $i>j$, maximize
lexicographically the tuple
\[
 (w,y,i,j),\qquad
 w=\|c_i-c_j\|_2^2
 \begin{cases}
 16,&c_i=\unif\text{ or }c_j=\unif,\\
 1,&\text{otherwise}.
 \end{cases}
\]
Replace vertex $j$ by the midpoint and recursively check that child; then
replace vertex $i$ by the midpoint and check the other child.  Both children
have depth one greater than their parent.
\end{enumerate}
Skipping Type P on cells having a uniform vertex is part of this specified
rule; no claim about the success or failure of that test on such cells is
needed.  Lemma~\ref{lem:subdivision} shows that successful termination proves
the required sign on the entire initial domain.

\subsection{The completed verification}

The verifier represents every coefficient matrix by integer numerators and
a common positive denominator, using arbitrary-precision integers for both.
Its polynomial arithmetic consists of convolution, degree elevation,
rational addition, and the exact substitutions of
Lemma~\ref{lem:subdivision}.  The fixed-width quantities used for indices,
counters, vertices, and small combinatorial factors remain within their integer
ranges.  Appendix~\ref{app:verifiers} proves these bounds and identifies
the implementation routines with the mathematical certificates.  Thus a
successful run verifies exact coefficient signs, without a numerical
tolerance.


\begin{proposition}[Finite sign verification]\label{prop:finite-verification}
For each row of \eqref{eq:rational-rows}, the decision rule above terminates
without failure, with the counts in Table~\ref{tab:verification-counts}.
Every terminal cell is accepted by Type T or Type P.  Consequently both
rational expressions \eqref{eq:rational-target} are nonnegative on
$\mathcal S\times\mathcal S$.
\end{proposition}

\begin{table}[ht]
\centering
\begin{tabular}{c|r|r|r|r|r}
Row & Cells visited & Maximum depth & Type T leaves & Type P leaves & Unresolved\\
\hline
(a)&4479&19&166&2074&0\\
(b)&1935&20&97&871&0
\end{tabular}
\caption{Exact subdivision counts.  The initial cell has depth zero;
visited cells include both internal nodes and terminal cells.}
\label{tab:verification-counts}
\end{table}

\begin{proof}
Run the integer verifier in Appendix~\ref{app:verifiers}, once in each
mode, with assertions enabled.  It returns the two lines
\begin{verbatim}
done 4479 19 lp=2074 lt=166
done 1935 20 lp=871 lt=97
\end{verbatim}
Here \texttt{lp} and \texttt{lt} count the terminal cells accepted by
Types P and T.  The substitutions are those of Lemma~\ref{lem:subdivision},
and Appendix~\ref{app:implementation} verifies their exact implementation
and the coefficient arithmetic.  Lemmas~\ref{lem:type-t} and~\ref{lem:type-p} certify the sign on
each accepted cell, and Lemma~\ref{lem:subdivision} proves coverage.  The
continuous extension at zero denominators established in
Section~\ref{sec:polynomial-certificates} includes the cell boundaries.
\end{proof}

For reproduction, the paper directory contains \path{verification/checks/certificate.cpp}
and \path{verification/checks/small_polynomials.py}, printed in full in
Appendix~\ref{app:verifiers}.  The first requires a C++17 compiler with
GCC-compatible standard-library headers and the Boost multiprecision
headers; the second uses only the Python 3 standard
library.  For example, from the paper's \path{verification/} directory run
\begin{verbatim}
mkdir -p artifacts
g++ -std=c++17 -O3 -UNDEBUG checks/certificate.cpp \
    -o artifacts/certificate
./artifacts/certificate
./artifacts/certificate b
python3 checks/small_polynomials.py
\end{verbatim}
The C++ program writes progress and its final result to standard error.
Its no-argument mode checks row (a); any argument selects row (b).
The Python program independently computes the four scalar coefficient
minima of Lemma~\ref{lem:scalar-checks}, using rational convolution, and
checks all 24 coordinate actions.  Assertions must remain enabled in both
programs: do not define \texttt{NDEBUG} or run Python with \texttt{-O}.
Together with the scalar checks, Proposition~\ref{prop:finite-verification}
completes the proof of Proposition~\ref{prop:polynomial-inequalities}.
